% SPDX-License-Identifier: CC-BY-4.0
\documentclass[11pt,letterpaper]{article}

\usepackage[T1]{fontenc}
\usepackage{amsmath,amssymb}
\usepackage{tikz}
\usepackage{pgfplots}
\pgfplotsset{compat=1.18}
\usepackage{algorithm}
\usepackage{algpseudocode}
\usepackage{booktabs}
\usepackage{enumitem}
\usepackage[hidelinks]{hyperref}

\newcommand{\N}{\mathbb{N}}
\newcommand{\Z}{\mathbb{Z}}
\newcommand{\Q}{\mathbb{Q}}
\newcommand{\R}{\mathbb{R}}
\newcommand{\C}{\mathbb{C}}
\newcommand{\Map}[3]{#1 \colon #2 \to #3}
\newcommand{\Mod}[1]{#1\text{-}\mathrm{Mod}}
\newcommand{\GroupHom}{\mathrm{Hom}_{\mathbf{Grp}}}
\newcommand{\RingHom}{\mathrm{Hom}_{\mathbf{Ring}}}

\newtheorem{theorem}{Theorem}
\newtheorem{lemma}[theorem]{Lemma}
\newtheorem{corollary}[theorem]{Corollary}
\newtheorem{definition}[theorem]{Definition}
\newtheorem{remark}[theorem]{Remark}
\newenvironment{proof}{\noindent\textit{Proof.}}{\hfill\(\square\)}

\title{Generalized Collatz Library}
\author{}
\date{\today}

\begin{document}
\maketitle
\section*{Full-orbit compression recovery}

\begin{abstract}
Removing a common power of the signed divisor gives an exact variable-clock
identity. We recover every intermediate raw state, prove preservation of
boundedness, full-sequence magnitude escape and eventual periodicity, and
identify a sufficient condition for bidirectional target capture. Primitive
periods require an additional exclusion of returns inside division blocks.
That exclusion also yields a bijection of periodic cycle classes. All statements
hold on the signed integer state space; no independence or coprimality
assumption is used.
\end{abstract}

\section{Parameters and the exact block identity}

% The compressed coefficients are lowercase, matching the formal parameterization.
\begin{definition}
Fix \(d,q,c\in\Z\), \(|d|\ge2\), and \(a\in\N\). For integers \(u,v\), define
\[
F_{u,v,d}(x)=\begin{cases}x/d,&d\mid x,\\ux+v,&d\nmid x.\end{cases}
\]
Put \(F=F_{d^a q,d^a c,d}\) and \(G=F_{q,c,d}\). Division precedes the affine
branch, including at zero, and retains the sign of \(d\). Define
\[
\tau(x)=\begin{cases}1,&d\mid x,\\a+1,&d\nmid x.\end{cases}
\]
\[
R(k,n)=\#\{j<k:d\nmid G^j(n)\},\qquad
T(k,n)=\sum_{j<k}\tau(G^j(n)).
\]
Empty sums and counts are zero.
\end{definition}

The common-power block identity \cite{AscendCompressionI} is
\begin{equation}\label{BlockIdentity}
F^{\tau(x)}(x)=G(x),\qquad
G^k(n)=F^{T(k,n)}(n),\qquad T(k,n)=k+aR(k,n).
\end{equation}
Indeed, division is unchanged. On an affine input, the first raw state is
\(d^a(qx+c)\); the following \(a\) instructions divide away the displayed
factors. Composing these blocks proves the iterate identity by induction.
This identity is not a one-step conjugacy.

\section{Every raw time and its intermediate state}

% The finite interval between consecutive accumulated times contains every raw time.
\begin{lemma}[Positive-clock coverage]\label{ClockCover}
For every \(n\in\Z\) and \(r\in\N\), there exist \(k,s\in\N\) such that
\[
k\le r,\qquad s<\tau(G^k(n)),\qquad r=T(k,n)+s.
\]
Moreover \(T(k,n)\) is strictly increasing in \(k\) and
\(k\le T(k,n)\le(a+1)k\). If \(r\ge T(K,n)\), any such block index satisfies
\(k\ge K\).
\end{lemma}
\begin{proof}
Each increment \(T(k+1,n)-T(k,n)=\tau(G^k(n))\) lies between one and \(a+1\).
Summing these inequalities gives both bounds and strict increase. For coverage,
induct on \(r\). At zero choose \(k=s=0\). Given a block containing \(r\),
if \(s+1<\tau(G^k(n))\), keep \(k\) and replace \(s\) by \(s+1\).
Otherwise \(s+1=\tau(G^k(n))\); use block \(k+1\) with offset zero.
The bound \(k\le r\) is preserved because each new block consumes at least
one raw step. Finally, if \(k<K\), then \(k+1\le K\), and the interior
condition implies \(r<T(k+1,n)\le T(K,n)\), contradicting \(r\ge T(K,n)\).
\end{proof}

% Exact signed quotients expose all hidden states, not just the endpoint.
\begin{lemma}[Division prefixes and affine intermediates]\label{Intermediates}
For \(b,s\in\N\) with \(s\le b\),
\[
F^s(d^b x)=d^{b-s}x.
\]
If \(d\nmid x\) and \(1\le s\le a+1\), then
\[
F^s(x)=d^{a+1-s}G(x).
\]
Consequently every positive-time interior of an affine block is divisible by
\(d\).
\end{lemma}
\begin{proof}
During the first \(s\) instructions on \(d^b x\), each current state contains
an explicit factor \(d\); exact signed division removes that factor. Induction
proves the first formula, including zero and negative \(d\).
For the second, the initial affine instruction gives \(d^aG(x)\).
Apply the first formula to the next \(s-1\) instructions. At an interior time
\(0<s<a+1\), the remaining exponent \(a+1-s\) is positive, so the state is
divisible by \(d\).
\end{proof}

% Both endpoints are retained to avoid losing a lower bound near a division run.
\begin{lemma}[Adjacent-state magnitude bounds]\label{MagnitudeBounds}
For every \(x\in\Z\) and \(0\le s<\tau(x)\),
\[
\min(|x|,|G(x)|)\le |F^s(x)|
\le |d|^a\max(|x|,|G(x)|).
\]
\end{lemma}
\begin{proof}
At \(s=0\), the state is \(x\), and \(|d|^a\ge1\). If \(s>0\), a
division block of length one is impossible. The block is affine and
Lemma~\ref{Intermediates} gives
\(|F^s(x)|=|d|^{a+1-s}|G(x)|\). Its factor lies between one and \(|d|^a\),
which proves both bounds. No inequality between \(|x|\) and \(|G(x)|\) is assumed.
\end{proof}

\section{Full-sequence orbit properties}

% The definitions specify the exact quantifiers used by the Lean statements.
\begin{definition}
For \(f:\Z\to\Z\), the orbit of \(n\) is magnitude bounded if
\(\exists M\in\N\ \forall r\in\N:\ |f^r(n)|\le M\). It escapes in magnitude
if \(\forall M\in\N\ \exists R\in\N\ \forall r\ge R:\ |f^r(n)|>M\).
It is eventually periodic if there exist \(u,h\in\N\), \(h>0\), with
\(f^{u+h}(n)=f^u(n)\).
\end{definition}

% Full recovery uses clock coverage rather than an assertion about sampled limits.
\begin{theorem}[QCD.TimeCompression.002.1]\label{Outcomes}
Fix \(n\in\Z\). Each of magnitude boundedness, full-sequence magnitude
escape and eventual periodicity holds for \(F\) if and only if that same
property holds for \(G\).
\end{theorem}
\begin{proof}
If the raw orbit is bounded by \(M\), its sampled compressed states are
bounded by \(M\) using \eqref{BlockIdentity}. Conversely, if all compressed
states are bounded by \(M\), cover any raw time by Lemma~\ref{ClockCover}.
The raw state is \(F^s(G^k(n))\); both adjacent compressed states have magnitude
at most \(M\). Lemma~\ref{MagnitudeBounds} bounds it by \(|d|^aM\).

For escape from \(F\) to \(G\), a raw cutoff \(R\) is also a compressed
cutoff: for \(k\ge R\), \(T(k,n)\ge k\ge R\). For the reverse direction,
fix \(M\) and choose \(K\) after which every compressed magnitude exceeds
\(M\). If \(r\ge T(K,n)\), its covering block has \(k\ge K\).
Both \(|G^k(n)|\) and \(|G^{k+1}(n)|\) exceed \(M\), so the lower bound in
Lemma~\ref{MagnitudeBounds} proves \(|F^r(n)|>M\).

For any deterministic map on \(\Z\), eventual periodicity implies boundedness:
a return after \(u\) steps with period \(h>0\) reduces every later iterate to
one of the finite first \(u+h\) iterates. Conversely, a magnitude bound \(M\)
places every state in the \(2M+1\) element interval \([-M,M]\).
Among the first \(2M+2\) states, two coincide, giving a positive return after
a finite prefix. Thus boundedness is equivalent to eventual periodicity;
its already proved transfer yields the third assertion.
\end{proof}

% Guarded closure is a sufficient target condition, not an assertion of necessity.
\begin{theorem}[QCD.TimeCompression.002.2]\label{Capture}
Let \(S\subseteq\Z\) satisfy
\[
x\in S,\ d\mid x\quad\Longrightarrow\quad x/d\in S.
\]
Then for each \(n\in\Z\), \(\exists r:F^r(n)\in S\) if and only if
\(\exists k:G^k(n)\in S\). Without this condition, compressed capture still
implies raw capture.
\end{theorem}
\begin{proof}
A compressed hit is a raw hit at its accumulated time. For a raw hit, cover its
time as \(T(k,n)+s\). At \(s=0\) it is the compressed state \(G^k(n)\).
At \(s>0\), Lemma~\ref{Intermediates} writes it as
\(d^{a+1-s}G^{k+1}(n)\). Repeatedly apply the target condition along those
remaining forced divisions to obtain \(G^{k+1}(n)\in S\).
\end{proof}

For \(d=2,q=3,c=1\), both \(S=\{1\}\) and \(S=\{1,2\}\) satisfy the
condition: one is not divisible by two, and two divides to one.
Thus the scaled family has the same exact target-one capture question as the
classical raw map, for each signed seed. No assertion that all positive seeds
are captured follows. An arbitrary converse is false: \(F_{4,4,2}(1)=8\), but
the \(F_{1,1,2}\) orbit of one remains in \(\{1,2\}\).

\section{Primitive periods}

% A return is distinguished from its least positive period throughout.
\begin{definition}
A primitive period of \(x\) for a map \(f\) is an integer \(h>0\) such that
\(f^h(x)=x\) and every positive return time is at least \(h\).
A periodic point is one admitting any positive return.
\end{definition}

% This elementary cancellation is valid only on points with a common return.
\begin{lemma}[Periodic predecessor uniqueness]\label{Predecessors}
If \(f^h(x)=x\), \(f^h(y)=y\), \(h>0\), and \(f(x)=f(y)\), then \(x=y\).
Points with separate positive return periods \(h_x,h_y\) have the common return
period \(h_xh_y\).
\end{lemma}
\begin{proof}
Apply \(f^{h-1}\) to the image equality. The two sides become \(x\) and \(y\).
A return propagates to every phase because iterates commute. Repeating the
\(h_x\)-return \(h_y\) times, and the \(h_y\)-return \(h_x\) times, gives
the common product period. No global injectivity of \(f\) is required.
\end{proof}

% The product period handles different cycles, which is needed for quotient injectivity.
\begin{lemma}[No periodic interior points]\label{NoInterior}
If \(x,y\) are periodic for \(G\), no positive-time interior of any expanded
block starting on the \(G\)-orbit of \(x\) equals \(y\).
\end{lemma}
\begin{proof}
First let \(z,y\) be periodic for \(G\) with \(d\nmid z\).
Suppose \(y=d^eG(z)\) for some \(e>0\). The first \(e-1\) instructions
from \(y\) are forced divisions and reach \(w=dG(z)\).
This phase is periodic with every return period of \(y\), and
\(G(w)=G(z)\). Choose a common positive return period for \(w,z\).
Lemma~\ref{Predecessors} gives \(w=z\), contrary to \(d\mid w\) and
\(d\nmid z\).

Now take an interior time of a block based at \(z=G^k(x)\).
It forces \(d\nmid z\), and Lemma~\ref{Intermediates} expresses the state
as \(d^eG(z)\) for \(e>0\). The preceding argument excludes every periodic
\(y\), even if its cycle is different from that of \(x\).
\end{proof}

% Minimality must also exclude shorter returns between sampled times.
\begin{theorem}[QCD.TimeCompression.002.3]\label{Primitive}
For \(n\in\Z\) and \(h\in\N\), \(h\) is a primitive period for
\(G\) at \(n\) if and only if \(h+aR(h,n)\) is a primitive period for
\(F\) at that same seed.
\end{theorem}
\begin{proof}
Suppose first that \(h\) is primitive for \(G\). Its expanded time
\(T(h,n)\) is positive and returns under \(F\).
Any other positive raw return is covered by a block. Lemma~\ref{NoInterior},
with the periodic target \(n\), forces its offset to be zero, so the return is
\(T(k,n)\) for a positive compressed return index \(k\). Primitiveness gives
\(h\le k\), and the increasing clock gives \(T(h,n)\le T(k,n)\).
Thus the expanded return is primitive.

Conversely, suppose \(T(h,n)\) is a primitive raw period. It is positive, so
\(h>0\), and \eqref{BlockIdentity} gives \(G^h(n)=n\).
Every positive compressed return \(k\) expands to a positive raw return
\(T(k,n)\), necessarily at least \(T(h,n)\). Strict increase of the clock
then gives \(h\le k\). The formula for \(T\) proves the displayed assertion.
\end{proof}

% Affine-phase recovery avoids the false converse at an arbitrary raw periodic seed.
\begin{theorem}[QCD.TimeCompression.002.4]\label{RawRecovery}
Every nonzero primitive raw cycle has a phase \(x=F^k(n)\), \(k<P\), with
\(d\nmid x\), and a primitive compressed period \(h\) at that phase, such that
\(P=h+aR(h,x)\). Zero has primitive period one for both maps.
\end{theorem}
\begin{proof}
If all \(P>0\) instructions of a return were divisions, their exact product
would give \(d^Pn=n\). Since \(|d|^P>1\), this forces \(n=0\).
A nonzero cycle therefore has an affine phase \(x\). The raw primitive period
is unchanged by that phase shift: a shorter return at \(x\), shifted back by
\(P-k\) instructions, would be a shorter return at \(n\).

Cover the positive raw return from \(x\). Any positive block interior is
divisible by \(d\), whereas \(x\) is not. Hence the return is a boundary
\(T(h,x)=P\), with \(h>0\) and \(G^h(x)=x\).
The reflection direction of Theorem~\ref{Primitive} proves compressed
primitiveness. Zero always takes division and remains zero; its least positive
return is one.
\end{proof}

For clarity, raw periodicity need not imply compressed periodicity at the
original seed. The raw map \(F_{4,4,2}\) has cycle
\(8\to4\to2\to1\to8\), but \(F_{1,1,2}\) sends eight through
\(8,4,2,1,2,1,\ldots\); eight is transient. The affine phase one supplies
the compressed cycle \(1\leftrightarrow2\), with raw period \(2+2\cdot1=4\).
A signed-divisor control is \(F_{2,2,-2}:1\to4\to-2\to1\), compressed to
\(F_{-1,-1,-2}:1\to-2\to1\), with period \(2+1\cdot1=3\).
For the homogeneous control \(F_{2,0,2}\), the raw cycle \(1\to2\to1\)
compresses to the fixed point one of \(F_{1,0,2}\); its expanded period is
\(1+1\cdot1=2\). In each example the compressed cycle class maps to the
specified raw cycle class. The zero class, whose period is one and affine
count is zero, is retained separately.

\section{Periodic cycle classes}

% Restriction to periodic points prevents transient singleton classes from being called cycles.
\begin{definition}
For a map \(f\), write \(x\sim_f y\) when there are nonnegative integers
\(r,s\) with \(f^r(x)=y\) and \(f^s(y)=x\). On the periodic point set,
this is an equivalence relation: zero witnesses reflexivity, exchanging paths
proves symmetry, and concatenating both directed paths proves transitivity.
Its quotient is the set of periodic cycle classes.
\end{definition}

% Both forward paths must expand, and reflection requires both representatives periodic.
\begin{theorem}[QCD.TimeCompression.002.5]\label{CycleClasses}
Expansion induces a bijection from the periodic cycle classes of \(G\) to
those of \(F\). It sends the class of a compressed-periodic point \(x\) to
its raw cycle class and includes the zero class.
\end{theorem}
\begin{proof}
A positive compressed return expands to a positive raw return. Also, any
compressed path to an endpoint expands with its own accumulated clock to a
raw path with that same endpoint. Thus expansion preserves both directions
of mutual reachability and defines a map on the quotient classes.

For injectivity, let \(x,y\) be periodic for \(G\), and suppose a raw path
from \(x\) reaches \(y\). Cover its time by a compressed block. The target
\(y\) is periodic, so Lemma~\ref{NoInterior} forces offset zero; the path is
already a compressed path. Apply the same reasoning to the reverse path.
Thus two compressed classes with one raw image were already equal.

For surjectivity, take any raw periodic point. If it is zero, use zero.
Otherwise the all-division calculation in Theorem~\ref{RawRecovery} supplies
an affine phase on its raw cycle. Its raw return is a compressed boundary,
with positive compressed index, so this phase is periodic for \(G\).
The original raw point reaches that phase within the cycle, and the phase
reaches it by the remaining instructions of the return. Their raw classes
coincide, giving a preimage for every raw class.
\end{proof}

\section{Reduction resource and exact finite certificates}

The original and replacement obligations have the same signed seeds and the
same divisor. The preceding theorems recover boundedness, magnitude escape and
eventual periodicity in both directions, capture under the sufficient target
condition, and primitive cycles after the stated phase selection. The height
\(H(u,v)=\max(|u|,|v|)\) satisfies
\[
H(d^a q,d^a c)=|d|^aH(q,c).
\]
For \(a>0\) and \((q,c)\ne(0,0)\), the right factor exceeds one and the
compressed height is positive, so removing the power strictly lowers this
natural-number resource. The cases \(a=0\) and \((q,c)=(0,0)\) have no strict
decrease. The finite block formulas and bounds above account for every skipped
raw time; no unresolved recovery oracle is used. The corresponding target
problem for \(G\) remains the replacement obligation, not an assumed theorem.

The exact Go checker uses signed arbitrary-precision integers and eight
independent seed chunks. It examines
\[
\begin{aligned}
d&\in\{-5,-4,-3,-2,2,3,4,5\},& a&\in\{0,1,2,3\},\\
q,c&\in\{-3,-2,-1,0,1,2,3\},& -30&\le n\le30.
\end{aligned}
\]
Each compressed search stops after 32 steps. The first repeated state selects
a primitive compressed cycle. If its length is \(h\) with \(r\) affine
instructions, the checker replays all \(h+ar\) raw steps, rejects any earlier
positive return to its selected phase, and checks the final endpoint.
The grid contains 95,648 seed-parameter cases; the reference run gives 58,864
expansion certificates, not 58,864 distinct cycles. No-repeat cases are
undecided, not certified escape. The finite checks illustrate the universally
proved statements; they do not replace their proofs or solve the remaining
capture problem.

The formal statements are organized as follows.
\begin{itemize}
\item \texttt{CompressionOutcomes.lean}: full-time recovery and capture.
\item \texttt{CompressionPeriods.lean}: primitive periods and affine-phase recovery.
\item \texttt{CompressionCycleClasses.lean}: periodic cycle quotients.
\end{itemize}
The corresponding main declarations include:
\begin{itemize}
\item \texttt{CompressionClassicalCaptureOneIff}, for the exact target one;
\item \texttt{CompressionPrimitivePeriodFormula}, for the expanded least period;
\item \texttt{CompressionCycleMapBijective}, for the cycle-class bijection.
\end{itemize}

\bibliographystyle{amsalpha}
\bibliography{CompressionRecovery}

\end{document}
